% =====================================================================
% Apriori Algorithm - Academic Report
% =====================================================================
\documentclass[11pt,a4paper]{article}

\usepackage[utf8]{inputenc}
\usepackage[T1]{fontenc}
\usepackage[english]{babel}
\usepackage{geometry}
\geometry{margin=2.4cm}
\usepackage{xcolor}
\usepackage{titlesec}
\usepackage{listings}
\usepackage{hyperref}
\usepackage{enumitem}
\usepackage{array}
\usepackage{booktabs}
\usepackage{fancyhdr}

\definecolor{EarthYellow}{HTML}{E1A95F}
\definecolor{DelftBlue}{HTML}{1F305E}
\definecolor{SoftYellow}{HTML}{F4E1C1}
\definecolor{CodeBg}{HTML}{FBF6EB}

\hypersetup{
  colorlinks=true,
  linkcolor=DelftBlue,
  urlcolor=DelftBlue,
  citecolor=DelftBlue
}

\titleformat{\section}
  {\color{DelftBlue}\Large\bfseries}{\thesection.}{0.6em}{}
  [\vspace{-4pt}\color{EarthYellow}\rule{\linewidth}{1.2pt}]

\titleformat{\subsection}
  {\color{DelftBlue}\large\bfseries}{\thesubsection}{0.5em}{}

\pagestyle{fancy}
\fancyhf{}
\renewcommand{\headrulewidth}{0.6pt}
\renewcommand{\headrule}{\hbox to\headwidth{\color{EarthYellow}\leaders\hrule height \headrulewidth\hfill}}
\fancyhead[L]{\textcolor{DelftBlue}{\small Apriori Data Mining Project}}
\fancyhead[R]{\textcolor{DelftBlue}{\small Ghoulam \& Ounes}}
\fancyfoot[C]{\textcolor{DelftBlue}{\thepage}}

\lstdefinestyle{cppstyle}{
  backgroundcolor=\color{CodeBg},
  basicstyle=\ttfamily\footnotesize\color{DelftBlue},
  keywordstyle=\color{DelftBlue}\bfseries,
  commentstyle=\color{EarthYellow}\itshape,
  stringstyle=\color{EarthYellow},
  showstringspaces=false,
  breaklines=true,
  numbers=left,
  numberstyle=\tiny\color{DelftBlue},
  frame=single,
  rulecolor=\color{EarthYellow},
  tabsize=2,
  captionpos=b,
  language=C++
}
\lstset{style=cppstyle}

% =====================================================================
\begin{document}

% ==================== TITLE PAGE ====================
\begin{titlepage}
\centering
\vspace*{1.5cm}

{\color{EarthYellow}\rule{\linewidth}{2pt}}
\vspace{0.4cm}

{\Huge\bfseries\color{DelftBlue} Apriori Algorithm \par}
\vspace{0.3cm}
{\LARGE\color{DelftBlue} Frequent Itemset Mining\par}
\vspace{0.3cm}

{\color{EarthYellow}\rule{\linewidth}{2pt}}

\vspace{2cm}
{\Large\itshape\color{DelftBlue} A Data Mining Academic Project\par}
\vspace{0.4cm}
{\large Implemented in C++ with a Next.js User Interface\par}

\vfill

\begin{tabular}{>{\color{DelftBlue}\bfseries}r l}
Authors : & Ghoulam Mohamed Said \\
          & Ounes Abdelfattah Tahar \\[6pt]
Teacher : & Talbi Omar \\[6pt]
Module  : & Data Mining \\[6pt]
Year    : & 2025 / 2026 \\
\end{tabular}

\vspace{1.5cm}
{\color{EarthYellow}\rule{0.5\linewidth}{1pt}}
\end{titlepage}

% ==================== ABSTRACT ====================
\section*{\color{DelftBlue}Abstract}
\vspace{-6pt}\textcolor{EarthYellow}{\rule{\linewidth}{1pt}}\vspace{6pt}

This report presents a complete implementation of the \textbf{Apriori
algorithm}, one of the foundational algorithms in the field of data mining
and association rule learning. The algorithm is implemented from scratch in
C++ without any external library, exposed to a web client through a small
Python (Flask) API bridge, and consumed by a clean Next.js front-end. The
project demonstrates how frequent itemsets are extracted from a dataset of
transactions using support counting, candidate generation, and the
\emph{Apriori property} for pruning.

% ==================== 1. INTRODUCTION ====================
\section{Introduction}
Data mining is the process of discovering useful patterns, correlations and
trends in large datasets. One of the most studied tasks in this field is
\textbf{association rule mining}, where the goal is to find items that
frequently appear together in transactions.

The \textbf{Apriori algorithm}, proposed by Agrawal and Srikant in 1994,
is the classical algorithm for this task. Its core observation is the
\emph{Apriori property}:

\begin{quote}
\itshape ``If an itemset is frequent, then all of its subsets must also be
frequent. Conversely, if a subset is infrequent, then any superset cannot
be frequent.''
\end{quote}

% ==================== 2. PROBLEM DEFINITION ====================
\section{Problem Definition}

Let $\mathcal{I} = \{i_1, i_2, \dots, i_m\}$ be a set of items and let
$\mathcal{T} = \{t_1, t_2, \dots, t_n\}$ be a set of transactions where
each $t_j \subseteq \mathcal{I}$.

\begin{itemize}[leftmargin=1.2em]
  \item The \textbf{support count} of $X$ is the number of transactions
        containing $X$.
  \item Given a \textbf{minimum support} $\sigma$, find every itemset $X$
        such that $\text{support}(X) \geq \sigma$.
\end{itemize}

% ==================== 3. METHODOLOGY ====================
\section{Methodology}

The algorithm proceeds level by level. $L_k$ denotes frequent
$k$-itemsets; $C_k$ denotes candidate $k$-itemsets.

\begin{enumerate}[leftmargin=1.4em]
  \item Scan the database once to count single items and build $L_1$.
  \item For $k \geq 2$:
    \begin{enumerate}
      \item Generate $C_k$ by joining $L_{k-1}$ with itself (join step).
      \item Remove every candidate with an infrequent $(k-1)$-subset
            (prune step).
      \item Scan the database to count support of surviving candidates.
      \item Keep only those with support $\geq \sigma$ to form $L_k$.
    \end{enumerate}
  \item Stop when $L_k = \emptyset$.
\end{enumerate}

% ==================== 4. ALGORITHM EXPLANATION ====================
\section{Algorithm Explanation (Step by Step)}

\subsection{Step 1 -- Build $L_1$}
Count each item across all transactions. Items meeting $\sigma$ form $L_1$.

\subsection{Step 2 -- Candidate generation}
Two itemsets in $L_{k-1}$ are joined if their first $k-2$ items match and
their last items differ (smaller first). This is the classical
\emph{apriori-gen} procedure.

\subsection{Step 3 -- Pruning}
For each candidate $c$ of size $k$, every $(k-1)$-subset must exist in
$L_{k-1}$. If any subset is missing, $c$ is discarded before any scan.

\subsection{Step 4 -- Support counting}
Scan the database and count transactions containing each surviving
candidate. Those meeting $\sigma$ become $L_k$.

\subsection{Step 5 -- Iterate}
Repeat with $k+1$ until no new frequent itemset is produced.

% ==================== 5. IMPLEMENTATION ====================
\section{Implementation Details}

\subsection{Project Architecture}
\begin{itemize}[leftmargin=1.4em]
  \item \textbf{C++ binary} (\texttt{apriori.cpp}): full algorithm,
        no external dependency, outputs JSON.
  \item \textbf{Flask bridge} (\texttt{backend/server.py}): receives HTTP
        requests, calls the binary as a subprocess, returns JSON.
  \item \textbf{Next.js frontend} (\texttt{frontend/}): upload or paste
        dataset, set minimum support, display results as table with
        support bars.
\end{itemize}

\subsection{Core C++ Functions}

\begin{lstlisting}[caption={Generate $L_1$ -- frequent 1-itemsets}]
vector<Itemset> generateL1(const vector<Transaction>& transactions,
                           int minSup, map<Itemset,int>& support) {
    map<string,int> itemCount;
    for (const auto& t : transactions)
        for (const auto& item : t)
            itemCount[item]++;
    vector<Itemset> L1;
    for (auto& kv : itemCount) {
        if (kv.second >= minSup) {
            Itemset s = { kv.first };
            L1.push_back(s);
            support[s] = kv.second;
        }
    }
    sort(L1.begin(), L1.end());
    return L1;
}
\end{lstlisting}

\begin{lstlisting}[caption={Candidate generation -- join step}]
vector<Itemset> aprioriGen(const vector<Itemset>& Lprev) {
    vector<Itemset> Ck;
    int n = Lprev.size();
    int k = Lprev.empty() ? 0 : (int)Lprev[0].size();
    for (int i = 0; i < n; i++)
      for (int j = i + 1; j < n; j++) {
        bool matchPrefix = true;
        for (int p = 0; p < k - 1; p++)
          if (Lprev[i][p] != Lprev[j][p]) { matchPrefix=false; break; }
        if (!matchPrefix) continue;
        if (Lprev[i][k-1] >= Lprev[j][k-1]) continue;
        Itemset c = Lprev[i];
        c.push_back(Lprev[j][k-1]);
        Ck.push_back(c);
      }
    return Ck;
}
\end{lstlisting}

\begin{lstlisting}[caption={Pruning -- Apriori property}]
vector<Itemset> pruneCandidates(const vector<Itemset>& Ck,
                                const vector<Itemset>& Lprev) {
    set<Itemset> Lset(Lprev.begin(), Lprev.end());
    vector<Itemset> pruned;
    for (const auto& c : Ck) {
        bool keep = true;
        for (size_t i = 0; i < c.size(); i++) {
            Itemset subset;
            for (size_t j = 0; j < c.size(); j++)
                if (j != i) subset.push_back(c[j]);
            if (Lset.find(subset)==Lset.end()) { keep=false; break; }
        }
        if (keep) pruned.push_back(c);
    }
    return pruned;
}
\end{lstlisting}

% ==================== 6. RESULTS ====================
\section{Results and Discussion}

\subsection{Example Dataset (8 transactions)}

\begin{center}
\begin{tabular}{|c|l|}
\hline
\textbf{TID} & \textbf{Items} \\
\hline
T1 & bread, milk, eggs \\
T2 & bread, butter \\
T3 & milk, eggs, butter \\
T4 & bread, milk, butter \\
T5 & bread, milk, eggs, butter \\
T6 & milk, eggs \\
T7 & bread, eggs \\
T8 & bread, milk \\
\hline
\end{tabular}
\end{center}

\subsection{Output -- minimum support $= 3$}

\paragraph{$L_1$:}
\begin{center}
\begin{tabular}{|l|c|}
\hline
\textbf{Itemset} & \textbf{Support} \\
\hline
\{bread\}  & 6 \\
\{milk\}   & 6 \\
\{eggs\}   & 5 \\
\{butter\} & 4 \\
\hline
\end{tabular}
\end{center}

\paragraph{$L_2$:}
\begin{center}
\begin{tabular}{|l|c|}
\hline
\textbf{Itemset} & \textbf{Support} \\
\hline
\{bread, milk\}   & 4 \\
\{eggs, milk\}    & 4 \\
\{bread, butter\} & 3 \\
\{bread, eggs\}   & 3 \\
\{butter, milk\}  & 3 \\
\hline
\end{tabular}
\end{center}

No 3-itemset reaches $\sigma = 3$; the algorithm stops at $L_2$.

\subsection{Discussion}
Pruning eliminates candidates early, reducing the number of database
scans. The algorithm scales linearly with the number of levels $k$, which
is its main bottleneck on very large datasets — motivating more advanced
methods such as FP-Growth.

% ==================== 7. CONCLUSION ====================
\section{Conclusion}

The Apriori algorithm was implemented from scratch in C++ and exposed
through a clean web interface using Next.js and a Flask bridge. The
implementation is short, modular, and easy to read, making it suitable for
educational purposes. Possible improvements include percentage-based
thresholds, association rule generation with confidence and lift, and
benchmarking against FP-Growth.

% ==================== REFERENCES ====================
\section*{\color{DelftBlue}References}
\vspace{-6pt}\textcolor{EarthYellow}{\rule{\linewidth}{1pt}}\vspace{6pt}

\begin{enumerate}[leftmargin=1.4em]
  \item R. Agrawal and R. Srikant, \emph{Fast Algorithms for Mining
        Association Rules}, VLDB, 1994.
  \item J. Han, M. Kamber and J. Pei, \emph{Data Mining: Concepts and
        Techniques}, 3rd ed., Morgan Kaufmann, 2011.
  \item P.-N. Tan, M. Steinbach and V. Kumar, \emph{Introduction to Data
        Mining}, Addison-Wesley, 2005.
  \item Next.js Documentation -- \url{https://nextjs.org/docs}
  \item Flask Documentation -- \url{https://flask.palletsprojects.com/}
\end{enumerate}

\end{document}
